\documentclass[conference]{IEEEtran}
\IEEEoverridecommandlockouts

% ---------- Paquetes ----------
\usepackage[spanish,es-tabla,es-noshorthands]{babel}
\usepackage[utf8]{inputenc}
\usepackage[T1]{fontenc}
\usepackage{cite}
\usepackage{amsmath,amssymb,amsfonts}
\usepackage{graphicx}
\usepackage{textcomp}
\usepackage{xcolor}
\usepackage{booktabs}
\usepackage{tabularx}
\usepackage{array}
\newcolumntype{Y}{>{\raggedright\arraybackslash}X}
\usepackage{url}
\usepackage{listings}
\usepackage{tikz}
\usetikzlibrary{positioning,arrows.meta,shapes.geometric,fit,calc}
\usepackage{hyperref}

\hypersetup{colorlinks=true,linkcolor=blue,filecolor=magenta,urlcolor=cyan,citecolor=blue}

\lstset{
  basicstyle=\ttfamily\scriptsize,
  breaklines=true,
  breakatwhitespace=false,
  columns=fullflexible,
  keepspaces=true,
  frame=single,
  framerule=0.2pt,
  xleftmargin=2pt,xrightmargin=2pt,
  commentstyle=\itshape\color{gray},
  showstringspaces=false
}

\tikzset{
  box/.style={draw,rounded corners=2pt,align=center,font=\scriptsize,minimum height=7mm,fill=blue!5,inner sep=2pt},
  dev/.style={draw,rounded corners=2pt,align=center,font=\scriptsize,minimum height=7mm,fill=gray!10,inner sep=2pt},
  srv/.style={draw,thick,rounded corners=2pt,align=center,font=\scriptsize,minimum height=7mm,fill=green!8,inner sep=2pt},
  dec/.style={draw,diamond,aspect=2.2,align=center,font=\scriptsize,fill=yellow!15,inner sep=1pt},
  lnk/.style={-{Stealth[length=2mm]},thin},
  lbl/.style={font=\tiny,midway,fill=white,inner sep=1pt}
}

\begin{document}

\title{Diseño e Implementación de Infraestructura de Red Local Autónoma para Atención Móvil Comunitaria en Entornos Rurales}


\author{\IEEEauthorblockN{Jose David Rodríguez Pinilla}
\IEEEauthorblockA{\textit{Código: A00407151} \\
\textit{Profesor: Andrés Navarro Cadavid, Ph.D.} \\
\textit{Infraestructura de Redes 2} \\
\textit{Departamento de Tecnologías de Información y Comunicaciones} \\
\textit{Universidad Icesi}\\
Cali, Colombia \\
jdrpjd@gmail.com}
}

\maketitle

\begin{abstract}
Las zonas rurales de difícil acceso suelen carecer de conectividad confiable a Internet y sufren interrupciones eléctricas. Este documento presenta el plan de ingeniería de una red local autónoma para una unidad móvil de atención comunitaria, sobre la cual opera una plataforma de atención desarrollada en Spring Boot y desplegada en contenedores. La solución se basa en un nodo único Raspberry Pi 5 (8~GB) que actúa como enrutador, servidor DHCP/DNS, portal cautivo, proxy/caché Squid y servidor de contenidos, con un SSD local de 2~TB ampliable a 4~TB. Se define una topología con segmentación L2/L3 (VLAN 802.1Q), un portal cautivo único que asigna perfiles de acceso (Comunidad, Educadores y Salud), una estrategia de caché con jerarquía de almacenamiento, políticas de refresco y manejo explícito de HTTPS, un mecanismo de conmutación online/offline, y un cronograma para un único integrante con ruta crítica, plan de compras escalonado y matriz de riesgos.
\end{abstract}

\begin{IEEEkeywords}
Redes rurales, portal cautivo, Squid, caché local, Raspberry Pi 5, Spring Boot, operación offline, redes comunitarias.
\end{IEEEkeywords}

%=====================================================================
\section{Introducción}
El acceso a servicios digitales en áreas rurales dispersas enfrenta ausencia de cobertura móvil de alta velocidad, altos costos de transmisión satelital e inestabilidad eléctrica. Una unidad móvil de atención comunitaria necesita, por tanto, una red propia que funcione con un enlace intermitente (modo \textit{online}) y en aislamiento total (modo \textit{offline}), ofreciendo a cada grupo de interés la experiencia de una red independiente con acceso a Internet cuando exista.

Este documento corresponde a la Fase~1 (Plan de Proyecto) y cubre: (i)~la arquitectura de red L2/L3 optimizada para el hardware elegido; (ii)~la estrategia de caché y almacenamiento; y (iii)~el cronograma, los recursos y la gestión de riesgos. Las Fases~2 y~3 construyen el prototipo y el sistema plenamente operativo. Cada decisión técnica se justifica con el contexto de operación: unidad móvil, energía limitada, enlace intermitente, un único integrante y presupuesto reducido.

\subsection{Alcance y requisitos}
La Tabla~\ref{tab:req} resume los requisitos derivados del enunciado y su atención en el diseño.

\begin{table}[htbp]
\caption{Requisitos y decisión de diseño asociada}
\label{tab:req}
\centering\footnotesize
\begin{tabularx}{\columnwidth}{l Y Y}
\toprule
\textbf{ID} & \textbf{Requisito} & \textbf{Decisión de diseño} \\
\midrule
RQ1 & Red independiente con acceso a Internet cuando exista enlace & RPi5 como gateway con NAT; WAN por USB-Ethernet (Sec.~\ref{sec:hw}) \\
RQ2 & Acceso diferenciado por grupo (Comunidad, Educadores, Salud) & Portal cautivo único con perfiles (Sec.~\ref{sec:portal}) \\
RQ3 & Proxy local sobre Mini PC o Raspberry Pi & Squid en RPi5 con caché jerárquica (Sec.~\ref{sec:cache}) \\
RQ4 & Almacenamiento masivo de 2--4~TB & SSD USB~3 de 2~TB (ampliable) + NVMe para SO \\
RQ5 & Operar sin Internet con el contenido ya almacenado & Conmutación online/offline (Sec.~\ref{sec:switch}) \\
RQ6 & Tolerar cortes eléctricos del vehículo & Batería auxiliar 12~V y apagado ordenado (Sec.~\ref{sec:power}) \\
RQ7 & Proteger datos de salud & SSID cifrado WPA3 y VLAN dedicada (Sec.~\ref{sec:l2}) \\
RQ8 & Operar la plataforma de atención de forma sincronizable & Spring Boot en contenedores, \textit{offline-first} (Sec.~\ref{sec:platform}) \\
\bottomrule
\end{tabularx}
\end{table}

%=====================================================================
\section{Arquitectura de Red}

\subsection{Selección de hardware: Raspberry Pi 5}\label{sec:hw}
Se selecciona la \textbf{Raspberry Pi 5 de 8~GB}~\cite{rpi5} como nodo único. La Tabla~\ref{tab:hw} la compara con la alternativa Mini PC (Intel Core i5) y justifica la decisión.

\begin{table}[htbp]
\caption{Comparación de plataformas (valores de referencia; precio de la RPi5 según el PVP oficial de 2026~\cite{rpiprice})}
\label{tab:hw}
\centering\footnotesize
\begin{tabularx}{\columnwidth}{l Y Y}
\toprule
\textbf{Criterio} & \textbf{Raspberry Pi 5 (8 GB)} & \textbf{Mini PC Core i5} \\
\midrule
CPU & ARM Cortex-A76, 4 núcleos, 2.4 GHz & x86\_64, 4--12 núcleos \\
Consumo típico & $\approx$4--12 W & $\approx$15--35 W \\
Ethernet & 1$\times$GbE nativo (+ USB 3 para la 2.ª) & 1--2$\times$GbE/2.5GbE nativos \\
Almacenamiento & PCIe 2.0 x1 (NVMe con HAT) + USB 3 & NVMe/SATA nativos \\
Video por hardware & Decodifica HEVC; sin H.264 por HW & Quick Sync (H.264/HEVC) \\
Costo aproximado & $\approx$US\$125 & $\approx$US\$250--400 \\
Replicabilidad & Alta (bajo costo y tamaño) & Media \\
\bottomrule
\end{tabularx}
\end{table}

\textbf{Justificación.} Una unidad móvil alimentada por una batería de 12~V privilegia el bajo consumo, el costo y la replicabilidad por comunidades; además, la compra escalonada de un único integrante favorece una inversión inicial baja. Las limitaciones de la RPi5 se compensan así:
\begin{itemize}
  \item \textit{Una sola interfaz GbE:} \texttt{eth0} transporta el \textit{trunk} 802.1Q (LAN) y un adaptador USB~3 Gigabit hace de WAN (\texttt{eth1}); el enlace LTE (decenas de Mbps) es mucho más lento que el puerto GbE, por lo que éste no es el cuello de botella.
  \item \textit{Sin decodificación H.264 por hardware:} el video se pre-transcodifica en lote (perfiles H.264 480p/720p) para que Jellyfin haga \textit{direct play} sin transcodificar en vivo (Fase~3).
  \item \textit{Almacenamiento:} SO, contenedores, base de datos y caché de objetos pequeños en NVMe (PCIe); contenido masivo en un \textbf{SSD USB~3 (SATA) de 2~TB}. Se elige SSD en lugar de disco mecánico porque la unidad es un vehículo en movimiento: los HDD son sensibles a vibración y golpes, consumen más y generan más calor. Un SSD portátil de bajo consumo ($\le$1~A) se alimenta del puerto USB de la RPi5 (alimentación de 5~V/5~A; ver Sec.~\ref{sec:power}).
  \item \textit{Térmica:} disipador activo oficial y carcasa ventilada por el calor del vehículo.
  \item \textit{Ajustes del SO para el almacenamiento:} \texttt{ext4} con \texttt{noatime}; reserva de \textit{root} reducida a 1\,\% en la partición de datos (\texttt{tune2fs -m 1}), pues el 5\,\% por defecto desperdiciaría $\approx$100~GB; \texttt{read\_ahead\_kb} elevado para lectura secuencial de video; \texttt{fstrim.timer} semanal si el puente USB soporta TRIM; y \texttt{zram} con \texttt{vm.swappiness} bajo para no escribir \textit{swap} en el NVMe.
\end{itemize}

\textbf{Criterio de migración.} Si en la Fase~3 la CPU supera de forma sostenida el 80\,\% bajo la carga objetivo, o se requiere transcodificación en vivo, el diseño se traslada a un Mini PC Core i5 sin cambios de arquitectura, porque todos los servicios corren sobre Debian y contenedores \texttt{multi-arch}.

\subsection{Topología física}\label{sec:topo}
La Fig.~\ref{fig:topo} muestra la topología en estrella. El nodo concentra los servicios; un \textit{switch} gestionable pequeño distribuye el \textit{trunk} hacia el punto de acceso (AP) y ofrece un puerto de administración.

\begin{figure*}[t]
\centering
\resizebox{\textwidth}{!}{%
\begin{tikzpicture}
 \node[dev] (net) at (0,0) {Internet\\(LTE/satélite)};
 \node[dev] (mod) at (3.1,0) {Módem /\\router LTE};
 \node[srv,minimum width=30mm] (pi) at (7,0) {\textbf{Raspberry Pi 5}\\eth1 (USB3): WAN\\eth0: trunk 802.1Q\\NVMe + SSD 2 TB};
 \node[box] (sw) at (11.2,0) {Switch\\gestionable};
 \node[box] (ap) at (14.2,1.5) {AP Wi-Fi 6\\SSID 1 + SSID 2};
 \node[dev] (adm) at (14.2,-1.5) {Laptop admin\\(VLAN 99)};
 \node[dev] (c1) at (17.2,2.3) {Celulares\\Comunidad/Educ.};
 \node[dev] (c2) at (17.2,0.7) {Tablets\\clínicas};
 \node[box] (ups) at (7,-2.2) {Batería 12 V + DC-DC\\(apagado ordenado)};
 \draw[lnk,<->] (net)--(mod);
 \draw[lnk,<->] (mod)--(pi) node[lbl,above,pos=.5,fill=none]{WAN};
 \draw[lnk,<->] (pi)--(sw) node[lbl,above,pos=.5,fill=none]{trunk V10/V30/V99};
 \draw[lnk,<->] (sw)--(ap) node[lbl,above left,pos=.5,fill=none]{trunk};
 \draw[lnk,<->] (sw)--(adm) node[lbl,below left,pos=.5,fill=none]{acceso V99};
 \draw[lnk,<->] (ap.east)--(c1.west) node[lbl,above,pos=.5]{SSID1 (V30)};
 \draw[lnk,<->] (ap.east)--(c2.west) node[lbl,below,pos=.5]{SSID2 (V10)};
 \draw[lnk] (ups)--(pi);
\end{tikzpicture}}
\caption{Topología física y lógica de la unidad móvil.}
\label{fig:topo}
\end{figure*}

\subsection{Diseño en Capa 2: segmentación y SSID}\label{sec:l2}
Se definen tres VLAN 802.1Q~\cite{ieee8021q} y dos SSID:

\begin{itemize}
  \item \textbf{VLAN 30 -- Portal (SSID~1 \texttt{Comunidad-Rural}, abierto):} todos los usuarios generales se conectan aquí; el portal cautivo asigna el perfil (Comunidad, Educador o Salud) tras la autenticación.
  \item \textbf{VLAN 10 -- Salud segura (SSID~2 \texttt{Salud-Segura}, WPA3-SAE o WPA2-PSK):} para equipos clínicos del personal. No usa portal, porque el cifrado debe darse en capa~2 y un SSID abierto expone el tráfico en el aire.
  \item \textbf{VLAN 99 -- Gestión (cableada):} SSH, métricas y administración del AP y del \textit{switch}.
\end{itemize}

El AP se conecta al \textit{switch} por \textit{trunk} (VLAN 10, 30, 99 etiquetadas; VLAN 99 como gestión del propio AP). En la RPi5 se crean subinterfaces \texttt{eth0.10}, \texttt{eth0.30} y \texttt{eth0.99} (\textit{router-on-a-stick}). Se activa aislamiento de clientes (\textit{client isolation}) en el SSID~1.

\subsection{Diseño en Capa 3: direccionamiento}\label{sec:l3}
Se usa el bloque privado \texttt{10.10.0.0/16} (RFC~1918~\cite{rfc1918}) con VLSM (Tabla~\ref{tab:ip}). La VLAN~20 (\texttt{10.10.20.0/24}) queda reservada para crecimiento (por ejemplo, un segundo AP).

\begin{table}[htbp]
\caption{Direccionamiento IP y segmentación L3}
\label{tab:ip}
\centering\footnotesize
\begin{tabularx}{\columnwidth}{l l l Y}
\toprule
\textbf{Segmento} & \textbf{Subred} & \textbf{Gateway} & \textbf{DHCP (capacidad)} \\
\midrule
V10 Salud segura & 10.10.10.0/24 & 10.10.10.1 & .50--.200 (151); lease 12 h \\
V30 Portal & 10.10.30.0/23 & 10.10.30.1 & 10.10.30.50--10.10.31.250 ($\approx$457); lease 2 h \\
V99 Gestión & 10.10.99.0/28 & 10.10.99.1 & .5--.14 (10); estático para AP/switch \\
WAN & DHCP del ISP & dinámico & Según módem LTE/satelital \\
\bottomrule
\end{tabularx}
\end{table}

\textbf{Dimensionamiento.} El espacio de direcciones de la VLAN~30 admite $\approx$450 dispositivos distintos atendidos durante una jornada (lease de 2~h, que reutiliza direcciones). La capacidad de radio es menor: un AP Wi-Fi 6 sostiene de forma realista 50--80 clientes activos simultáneos, por lo que el objetivo de diseño es \textbf{$\le$60 clientes activos por AP}, ampliable con un segundo AP en la VLAN~20 reservada.

Los clientes usan como DNS y gateway la RPi5 en su subred. Sobre la WAN se aplica NAT (\textit{masquerade}). Las políticas de reenvío entre VLAN se implementan en \texttt{nftables}~\cite{nft}: V30$\to$V10 y V10$\to$V30 denegados; V99 sólo accesible desde la propia V99; las VLAN de usuario pueden llegar a los servicios locales (cadena \texttt{input}, de modo que no se rompe el aislamiento entre V10 y V30) y a la WAN (V30 sólo tras autenticarse en el portal). Además se bloquean DoT (TCP/853) y QUIC (UDP/443) para que los clientes recurran a DNS local y a TLS sobre TCP, único tráfico que Squid puede gestionar (Sec.~\ref{sec:https}).

\textbf{Capacidad de la ruta de datos.} Un flujo de video 720p H.264 pre-transcodificado ocupa $\approx$3~Mbps, de modo que 20 flujos simultáneos consumen $\approx$60~Mbps, cerca del 6\,\% del \textit{trunk} GbE ($\approx$940~Mbps útiles); el límite real es el radio Wi-Fi, no el puerto del nodo. El tráfico hacia Internet queda acotado por el enlace LTE, de ahí los límites por perfil de la Tabla~\ref{tab:roles}.

\subsection{Roles del servidor y de los clientes}
Los servicios de red críticos corren \textbf{nativos en el host} (acceso directo a interfaces y a \textit{netfilter}, y arranque independiente de Docker); los servicios de aplicación corren en \textbf{contenedores}.
\begin{itemize}
  \item \textbf{RPi5 (host):} router/NAT y firewall (\texttt{nftables}); \texttt{dnsmasq}~\cite{dnsmasq} (DHCP+DNS); \texttt{openNDS}~\cite{opennds} (portal cautivo); \texttt{Squid}~\cite{squid} (proxy/caché); \texttt{Nginx}~\cite{nginx} (\textit{reverse proxy} de \texttt{portal.internal}, \texttt{salud.internal}, \texttt{biblioteca.internal} y \texttt{video.internal}).
  \item \textbf{RPi5 (contenedores):} plataforma de atención en Spring Boot, PostgreSQL, Kolibri, Kiwix-serve y Jellyfin.
  \item \textbf{AP:} sólo capa~2 (SSID$\leftrightarrow$VLAN); no hace enrutamiento ni DHCP.
  \item \textbf{Clientes:} celulares, tablets y portátiles; sin software especial. Los de Salud usan la plataforma vía navegador por HTTPS.
\end{itemize}

\subsection{Plataforma de atención: Spring Boot en contenedores}\label{sec:platform}
La plataforma de atención del proyecto de Plataformas se implementa con \textbf{Spring Boot 3.x sobre Java 21}~\cite{springboot} y se despliega como imagen multi-arquitectura \texttt{linux/arm64} construida con Docker Buildx~\cite{buildx}, de modo que la misma imagen se desarrolla en una laptop x86 y corre en la RPi5. Las decisiones principales son:
\begin{itemize}
  \item \textbf{Base de datos:} PostgreSQL~\cite{postgres} en contenedor, con volumen en NVMe y memoria acotada (\texttt{shared\_buffers} $\approx$256~MB). Es suficiente para la concurrencia esperada y soporta transacciones para historias clínicas.
  \item \textbf{Orquestación:} Docker Compose~\cite{compose} en un solo nodo; Kubernetes añadiría consumo de RAM y complejidad sin beneficio en un único equipo.
  \item \textbf{Memoria de la JVM:} límite de contenedor de $\approx$1~GB con \texttt{-Xmx768m}, para convivir con Squid y Jellyfin en 8~GB (Tabla~\ref{tab:ram}).
  \item \textbf{Modo \textit{offline-first}:} la base local es la fuente de verdad durante la jornada. Los cambios se registran en una tabla \textit{outbox} con identificadores UUID y se sincronizan por REST idempotente con el servicio central cuando hay enlace; los registros clínicos son \textit{append-only} para evitar conflictos de escritura.
  \item \textbf{Seguridad y monitoreo:} Spring Security con roles, TLS terminado en Nginx y \textit{Actuator} (\textit{health}) integrado al monitoreo de la Sec.~\ref{sec:switch}.
\end{itemize}

\begin{table}[htbp]
\caption{Presupuesto de memoria RAM (8 GB, estimado)}
\label{tab:ram}
\centering\footnotesize
\begin{tabularx}{\columnwidth}{Y l l}
\toprule
\textbf{Servicio} & \textbf{Ejecución} & \textbf{RAM (GB)} \\
\midrule
SO + Docker & host & 0.8 \\
Squid (\texttt{cache\_mem} 1 GB + índice y proceso) & host & 1.5 \\
dnsmasq, openNDS, Nginx & host & 0.3 \\
Plataforma Spring Boot & contenedor & 1.0 \\
PostgreSQL & contenedor & 0.75 \\
Jellyfin (\textit{direct play}) & contenedor & 0.8 \\
Kolibri & contenedor & 0.7 \\
Kiwix-serve & contenedor & 0.3 \\
\textbf{Total comprometido} & & \textbf{6.15} \\
Libre (caché de páginas del SO) & & 1.85 \\
\bottomrule
\end{tabularx}
\end{table}

\subsection{Portal cautivo único con perfiles}\label{sec:portal}
El enunciado exige que el portal dé acceso a los distintos grupos de interés. Se evaluaron tres opciones (Tabla~\ref{tab:portal}).

\begin{table}[htbp]
\caption{Opciones de portal cautivo}
\label{tab:portal}
\centering\footnotesize
\begin{tabularx}{\columnwidth}{l Y Y Y}
\toprule
\textbf{Criterio} & \textbf{A: portal por SSID} & \textbf{B: portal único} & \textbf{C: único + SSID salud (elegida)} \\
\midrule
Cumple enunciado & Parcial & Sí & Sí \\
Experiencia de usuario & Confusa (3 redes) & Simple & Simple \\
Cifrado de datos clínicos & Sólo si SSID cifrado & No (SSID abierto) & Sí (SSID~2 WPA3) \\
Complejidad & Alta & Baja & Media \\
\bottomrule
\end{tabularx}
\end{table}

\textbf{Decisión: opción C.} Un único portal en el SSID~1 presenta tres entradas: \textit{Comunidad} (aceptar términos), \textit{Educador} y \textit{Salud} (usuario/clave locales). Tras autenticar, un \textit{script} BinAuth/FAS de openNDS inserta la IP del cliente en el conjunto \texttt{nftables} del perfil (\texttt{role\_com}, \texttt{role\_edu}, \texttt{role\_sal}) y aplica límites de tiempo y ancho de banda (Tabla~\ref{tab:roles}). El perfil Salud del portal es para personal que usa dispositivos propios (por ejemplo, promotores); su acceso a la plataforma es siempre por HTTPS, y los equipos clínicos del personal deben usar el SSID~2. Los perfiles se aplican en L3 porque una VLAN dinámica requeriría RADIUS y 802.1X, innecesarios para este alcance. El soporte de la opción DHCP~114 (RFC~8910~\cite{rfc8910}) para mejorar la detección en Android/iOS es optativo y se valida en pruebas.

\begin{table}[htbp]
\caption{Perfiles de acceso (parámetros ajustables)}
\label{tab:roles}
\centering\footnotesize
\begin{tabularx}{\columnwidth}{l l Y Y}
\toprule
\textbf{Perfil} & \textbf{Autent.} & \textbf{Acceso} & \textbf{Límites WAN} \\
\midrule
Comunidad & Términos & Biblioteca, educación abierta, caché & Sesión 2 h; 1--2 Mbps por cliente \\
Educador & Usuario/clave & Anterior + Kolibri y Jellyfin & Sesión 12 h; 5 Mbps \\
Salud (portal) & Usuario/clave & Plataforma de salud, sin límite local & Sesión 12 h; prioridad alta \\
\bottomrule
\end{tabularx}
\end{table}

Los límites aplican al tráfico hacia Internet con \texttt{tc}/HTB jerárquico: la clase padre es la capacidad medida del enlace LTE y las clases hijas son los perfiles. El tráfico hacia servicios locales no se estrangula, ya que no consume el enlace.

\subsection{Flujo de datos}
La Fig.~\ref{fig:flow} resume el recorrido de una petición.

\begin{figure*}[t]
\centering
\resizebox{\textwidth}{!}{%
\begin{tikzpicture}
 \node[dev] (cli) at (0,0) {Cliente\\(V30)};
 \node[dec] (auth) at (3.2,0) {¿Autenticado\\en openNDS?};
 \node[box] (port) at (3.2,2.2) {Portal cautivo\\(perfil: Com/Educ/Salud)};
 \node[dec] (dst) at (7.2,0) {¿Destino?};
 \node[srv] (ng) at (7.2,2.2) {Nginx\\servicios locales};
 \node[box] (https) at (7.2,-2.2) {Squid :3130\\(HTTPS \textit{peek+splice})};
 \node[box] (sq) at (10.7,0) {Squid :3129\\(HTTP interceptado)};
 \node[dec] (hit) at (14,0) {¿HIT en\\caché?};
 \node[box] (serve) at (17,0) {Sirve desde\\NVMe / SSD};
 \node[dec] (on) at (14,-2.2) {¿Modo\\online?};
 \node[box] (up) at (17,-2.2) {Descarga de\\Internet + guarda};
 \node[box] (err) at (14,-4.2) {Aviso local\\«sin conexión»};
 \draw[lnk] (cli)--(auth);
 \draw[lnk] (auth)--(port) node[lbl,pos=.5]{No};
 \draw[lnk] (auth)--(dst) node[lbl,pos=.5]{Sí};
 \draw[lnk] (dst)--(ng) node[lbl,pos=.5]{*.internal};
 \draw[lnk] (dst)--(sq) node[lbl,pos=.5]{HTTP};
 \draw[lnk,dashed] (dst)--(https) node[lbl,pos=.5]{443};
 \draw[lnk] (sq)--(hit);
 \draw[lnk] (hit)--(serve) node[lbl,pos=.5]{Sí};
 \draw[lnk] (hit)--(on) node[lbl,pos=.5]{No};
 \draw[lnk] (on)--(up) node[lbl,pos=.5]{Sí};
 \draw[lnk] (on)--(err) node[lbl,pos=.5]{No};
\end{tikzpicture}}
\caption{Flujo de datos de un cliente del SSID Comunidad-Rural.}
\label{fig:flow}
\end{figure*}

El DNS se resuelve en \texttt{dnsmasq}: los dominios internos devuelven la IP del nodo y los externos se reenvían al \textit{upstream} sólo en modo online. Los nombres internos usan el sufijo \texttt{.internal}, reservado para uso privado~\cite{internaltld}, y no \texttt{.local}, que el RFC~6762~\cite{rfc6762} destina a mDNS y que Android e iOS no siempre resuelven por DNS unicast. Un fragmento representativo de configuración:

\begin{lstlisting}
interface=eth0.30
interface=eth0.10
bind-dynamic
domain-needed
bogus-priv
local=/internal/
dhcp-range=set:portal,10.10.30.50,10.10.31.250,255.255.254.0,2h
dhcp-range=set:salud,10.10.10.50,10.10.10.200,255.255.255.0,12h
dhcp-option=tag:portal,option:router,10.10.30.1
dhcp-option=tag:portal,option:dns-server,10.10.30.1
dhcp-option=tag:salud,option:router,10.10.10.1
dhcp-option=tag:salud,option:dns-server,10.10.10.1
address=/portal.internal/10.10.30.1
address=/biblioteca.internal/10.10.30.1
address=/salud.internal/10.10.30.1
address=/video.internal/10.10.30.1
cache-size=10000
# upstream (online/offline) en archivo aparte
conf-file=/etc/dnsmasq.d/upstream.conf
\end{lstlisting}

%=====================================================================
\section{Estrategia de Caché y Almacenamiento}\label{sec:cache}

\subsection{Jerarquía de almacenamiento}
\begin{enumerate}
  \item \textbf{Nivel 1 -- RAM:} \texttt{cache\_mem} de Squid (1~GB) y caché DNS.
  \item \textbf{Nivel 2 -- NVMe (256--512~GB):} SO, contenedores, base de datos y \textit{cache\_dir} \texttt{rock} para objetos pequeños ($\le$256~KB).
  \item \textbf{Nivel 3 -- SSD USB (2~TB, ampliable a 4~TB):} \textit{cache\_dir} \texttt{aufs} (objetos grandes), repositorios estáticos (ZIM de Kiwix~\cite{kiwix}, cursos de Kolibri~\cite{kolibri}), videos de Jellyfin~\cite{jellyfin} y copias de respaldo.
\end{enumerate}

Cada área vive en una \textbf{partición separada} (LVM o particiones fijas, \texttt{ext4}, \texttt{noatime}, \texttt{data=ordered}) para que el crecimiento de la caché no llene el contenido curado. La Tabla~\ref{tab:budget} presenta el presupuesto de espacio. Un SSD de 2~TB nominales ofrece $\approx$1860~GiB útiles tras el formateo, base de la Tabla~\ref{tab:budget}. Se asigna una caché web moderada (200~GB en el SSD) y se prioriza el contenido curado: como la mayor parte del tráfico es HTTPS (Sec.~\ref{sec:https}), el valor del caché web es limitado, y además Squid mantiene en RAM un índice de aproximadamente 100~bytes por objeto, lo que en 8~GB impone acotar el tamaño de la caché.

\begin{table}[htbp]
\caption{Presupuesto de espacio, SSD de 2 TB (estimado; validar contra los contenidos reales)}
\label{tab:budget}
\centering\footnotesize
\begin{tabularx}{\columnwidth}{Y l}
\toprule
\textbf{Área} & \textbf{Asignación} \\
\midrule
Caché Squid (\texttt{aufs}) & 200 GB \\
Video educativo (Jellyfin, pre-transcodificado) & 1000 GB \\
Kiwix (Wikipedia en español y otros ZIM) & 100 GB \\
Kolibri (cursos y recursos) & 150 GB \\
Manuales y guías de salud (PDF) & 20 GB \\
Respaldos (BD, configuración) & 100 GB \\
Reserva libre ($\ge$15\,\% de $\approx$1860 GB útiles) & $\approx$290 GB \\
\bottomrule
\end{tabularx}
\end{table}

\subsection{Cuotas, purga y ampliación a 4~TB}
Las cuotas se imponen por partición (una por área de la Tabla~\ref{tab:budget}), de modo que ningún servicio consuma el espacio de otro. Un \texttt{systemd timer} diario ejecuta un \textit{script} que: (i)~elimina del caché de Jellyfin las miniaturas y temporales no accedidos en 30 días; (ii)~rota los registros de Squid y openNDS con retención de 30 días, en el NVMe; (iii)~conserva sólo las últimas 7 copias de respaldo de la BD; y (iv)~emite una alerta cuando una partición supera el 85\,\% (R7). La caché de Squid se purga sola mediante \texttt{LFUDA} y \texttt{cache\_swap}. Para llegar a 4~TB se añade un segundo SSD USB como biblioteca adicional de Jellyfin y Kolibri (admiten varias rutas), sin ensanchar volúmenes LVM entre discos, para que la falla de uno no invalide el otro.

\subsection{Políticas de refresco y configuración de Squid}
Squid opera con un puerto explícito (3128), uno interceptado para HTTP (3129) y uno interceptado para HTTPS (3130), hacia los que \texttt{nftables} redirige los puertos~80 y~443 de los clientes autenticados cuyo destino no es la red interna~\cite{squid}. Se instala el paquete \texttt{squid-openssl} de Debian, necesario para \textit{ssl-bump}~\cite{sslpeek}. Un núcleo de configuración:

\begin{lstlisting}
http_port 3128
http_port 3129 intercept
https_port 3130 intercept ssl-bump tls-cert=/etc/squid/splice.pem
acl step1 at_step SslBump1
ssl_bump peek step1
ssl_bump splice all
cache_mem 1024 MB
memory_replacement_policy heap GDSF
maximum_object_size 512 MB
cache_replacement_policy heap LFUDA
cache_dir rock /mnt/nvme/squid 40000 max-size=262144
cache_dir aufs /mnt/data/squid 200000 16 256 min-size=262145
cache_swap_low 90
cache_swap_high 95
refresh_pattern -i \.(pdf|epub|zim|mp4|webm|jpe?g|png|gif|css|js)$ 10080 90% 43200 override-expire reload-into-ims
refresh_pattern -i (/cgi-bin/|\?) 0 0% 0
refresh_pattern . 0 20% 4320
max_stale 30 days
include /etc/squid/mode.conf  # offline_mode on|off
\end{lstlisting}

La lógica de actualización es: (i)~el contenido estático se considera fresco 7 días como mínimo y hasta 30; (ii)~con enlace, los objetos vencidos se revalidan con el origen (IMS); (iii)~sin enlace, se sirven aunque estén vencidos (\texttt{offline\_mode on} y \texttt{max\_stale}); (iv)~el reemplazo usa \texttt{heap LFUDA}, que favorece objetos populares y grandes de alto valor; y (v)~el \textit{cache\_swap} 90/95\,\% inicia la purga antes de saturar la partición.

\subsection{Manejo de HTTPS}\label{sec:https}
La mayor parte del tráfico web actual es HTTPS, por lo que un caché sólo HTTP guardaría muy poco. Instalar una CA propia en celulares de la comunidad para hacer \textit{SSL-Bump} no es viable ni deseable (fricción, privacidad, riesgo en telemedicina). Se decide:
\begin{enumerate}
  \item \textbf{Túnel por defecto (\textit{peek and splice}):} el 443 de los clientes autenticados llega a Squid (puerto 3130), que lee el SNI del \textit{ClientHello}~\cite{sslpeek} para aplicar listas de dominios y registro, y luego deja pasar el flujo TLS sin descifrarlo. QUIC (UDP/443) se bloquea para forzar TLS sobre TCP.
  \item \textbf{Espejos locales} para el contenido pesado y crítico (Wikipedia/Kiwix, Kolibri, manuales de salud): se sirven por HTTPS desde \texttt{Nginx} con un certificado de la CA del proyecto, instalable de forma voluntaria; los dominios \texttt{.internal} no dependen de Internet. Los sitios estáticos educativos se espejan con \texttt{wget --mirror} o HTTrack y se publican por Nginx.
  \item \textbf{SSL-Bump opcional} sólo en equipos administrados del proyecto (educadores), nunca para el perfil Salud.
\end{enumerate}
El contenido HTTPS externo no navegado previamente no estará disponible offline; por eso la pre-descarga de contenido curado (siguiente sección) es el mecanismo principal.

\subsection{Descarga, persistencia y sincronización}
\textbf{Siembra inicial.} El contenido masivo ($\approx$1.3~TB) se carga una sola vez antes del despliegue, por Ethernet en la red del campus o del hogar, donde no hay límite de datos. \textbf{Actualización en campo.} El enlace LTE (plan de datos prepago de $\approx$30~GB al mes) se reserva para incrementales, con una cuota de sincronización de $\approx$1~GB/día configurable. Con enlace y en horas de baja demanda, un servicio programado (\texttt{systemd timers}) ejecuta: (a)~actualización de ZIM de Kiwix y de cursos de Kolibri; (b)~descarga de nuevos videos y su transcodificación en lote a H.264; (c)~sincronización con \texttt{rsync} de manuales, con reintentos reanudables; (d)~sincronización de la cola \textit{outbox} de la plataforma hacia el servicio central (Sec.~\ref{sec:platform}); y (e)~precarga (\textit{cache warming}) de la lista curada de sitios HTTP y documentos de salud, solicitados a través del propio proxy para poblar la caché.

\subsection{Conmutación online/offline}\label{sec:switch}
Un servicio de monitoreo comprueba la conectividad cada 30~s con tres señales: ping a \texttt{1.1.1.1}, resolución DNS y una petición HTTP a un endpoint de verificación. Con histéresis (3 fallos, $\approx$90~s, para pasar a offline; 2 éxitos, $\approx$60~s, para volver a online) se evitan oscilaciones. El cambio escribe \texttt{offline\_mode on|off} en un archivo incluido por \texttt{squid.conf}, ejecuta \texttt{squid -k reconfigure} y activa o detiene la cola de sincronización. En modo offline, \texttt{dnsmasq} deja de reenviar al \textit{upstream} para evitar \textit{timeouts}, y las páginas no disponibles muestran un aviso local en lugar de un error de red.

%=====================================================================
\section{Energía y Seguridad Básica}\label{sec:power}
\textbf{Presupuesto de potencia (estimado, peor caso):} RPi5 $\approx$12~W (incluye el adaptador USB), NVMe $\approx$2~W, SSD $\approx$3~W, AP $\approx$10~W, \textit{switch} $\approx$4~W y router LTE $\approx$8~W: total $\approx$39~W; el consumo típico se estima en $\approx$27~W. Se decide una \textbf{batería auxiliar LiFePO$_4$ de 12~V/20~Ah} ($\approx$256~Wh, $\approx$80\,\% utilizable $\approx$205~Wh y, con eficiencia de conversión de $\approx$90\,\%, $\approx$185~Wh), cargada por un cargador DC-DC desde el alternador del vehículo, y un conversor 12~V$\to$5~V/5~A para la RPi5; los demás equipos se alimentan a 12~V o con su propio DC-DC según su tensión de entrada, con fusible y protección contra transitorios en la entrada. Esto da $\approx$4.7~h de autonomía en el peor caso y $\approx$6.8~h en uso típico, con el motor apagado; en marcha el alternador repone la batería. Un monitor de voltaje inicia un apagado ordenado al llegar a un umbral de carga, protegiendo el sistema de archivos. Como este conversor no negocia USB-PD, la RPi5 limitaría los puertos USB a 600~mA, insuficiente para el SSD; se corrige con \texttt{usb\_max\_current\_enable=1} en \texttt{config.txt}~\cite{rpipower}. Los valores se confirmarán con medición en la Fase~3.

\textbf{Seguridad básica:} acceso SSH sólo por VLAN~99 y con llaves; contraseñas distintas por servicio; cuenta de administración sin acceso remoto como \textit{root} y con límite de intentos de SSH; actualizaciones automáticas de seguridad (\texttt{unattended-upgrades}) cuando haya enlace; firewall por defecto en \textit{drop}; HTTPS interno con certificados de la CA del proyecto; datos de salud con respaldo cifrado. Los registros del portal (MAC, IP, perfil) se conservan 30 días y el portal muestra el aviso de tratamiento de datos personales conforme a la Ley~1581 de 2012~\cite{ley1581}.

%=====================================================================
\section{Cronograma, Ruta Crítica y Recursos}

\subsection{Plan de trabajo}
Al ser un único integrante, todas las actividades se asignan a un solo recurso y se programan de forma secuencial (nivelación de recursos); las compras, descargas y transcodificaciones se ejecutan en segundo plano sin ocupar la semana de trabajo. La Tabla~\ref{tab:tasks} detalla las actividades y la Fig.~\ref{fig:gantt} el cronograma de 11 semanas.

\begin{table}[t]
\caption{Actividades, dependencias y entregables. Holgura lógica: sin restricción de recursos (compras disponibles a tiempo). Con un único recurso, toda la secuencia queda en la cadena crítica.}
\label{tab:tasks}
\centering\footnotesize
\begin{tabularx}{\textwidth}{l Y c l c c Y}
\toprule
\textbf{ID} & \textbf{Actividad} & \textbf{Dur. (sem)} & \textbf{Pred.} & \textbf{Semanas} & \textbf{Holgura lógica (sem)} & \textbf{Entregable / hito} \\
\midrule
T1 & Diseño y plan (informe, diagramas, riesgos, BOM) & 3 & -- & 1--3 & 0 & Informe Fase 1 \\
T2 & SO, particionado y montaje del SSD; Docker & 1 & T1, Lote 1 & 4 & 0 & Nodo arranca y monta 2 TB; lectura secuencial $\ge$300~MB/s (\texttt{fio}) \\
T3 & Red base: VLAN, nftables, dnsmasq, AP & 1 & T2, Lote 2 & 5 & 0 & Clientes con DHCP/DNS local en V10 y V30 \\
T4 & Squid + Nginx + caché de contenido estático & 1 & T3 & 6 & 0 & Sitios/PDF servidos offline; HIT verificado en \texttt{access.log} \\
T5 & Portal cautivo básico con perfiles (openNDS) & 1 & T3 & 7 & 1 & Portal funcional con 3 perfiles (fin Fase 2) \\
T6 & Plataforma Spring Boot en contenedores + PostgreSQL & 1 & T3 & 8 & 1 & Plataforma en \texttt{salud.internal} \\
T7 & Kiwix, Kolibri, Jellyfin y pre-transcodificación & 1 & T2, T4 & 9 & 0 & Contenido curado; video sin \textit{buffering} con 10 clientes \\
T8 & Dashboard categorizado con búsqueda local & 0.5 & T5, T7 & 10 & 0 & Portal de contenidos con búsqueda \\
T9 & Conmutación online/offline y sincronización & 0.5 & T4, T6, Lote 3 & 10 & 1 & Modo automático; conmuta en $\le$2~min \\
T10 & Pruebas de carga, autonomía, endurecimiento y entrega & 1 & T8, T9 & 11 & 0 & Pruebas superadas; sistema plenamente operativo \\
\bottomrule
\end{tabularx}
\end{table}

\begin{figure}[t]
\centering
\begin{tikzpicture}[x=0.52cm,y=-0.42cm,font=\scriptsize]
 \foreach \w in {1,...,11} {
   \node[font=\tiny] at (\w-0.5+4,-0.2) {S\w};
 }
 \foreach \w in {0,...,11} { \draw[gray!30] (\w+4,0.2) -- (\w+4,13.8); }
 \newcommand{\gbar}[5]{%
   \node[anchor=east] at (3.8,#1) {#4};
   \fill[#5] (#2+4,#1-0.3) rectangle (#3+4,#1+0.3);}
 \gbar{1}{0}{3}{T1 Diseño y plan}{red!70}
 \gbar{2}{3}{4}{T2 SO + almacenamiento}{red!70}
 \gbar{3}{4}{5}{T3 Red base L2/L3}{red!70}
 \gbar{4}{5}{6}{T4 Squid + Nginx}{red!70}
 \gbar{5}{6}{7}{T5 Portal cautivo}{red!70}
 \gbar{6}{7}{8}{T6 Spring Boot (contenedores)}{red!70}
 \gbar{7}{8}{9}{T7 Contenidos y multimedia}{red!70}
 \gbar{8}{9}{9.5}{T8 Dashboard y búsqueda}{red!70}
 \gbar{9}{9.5}{10}{T9 Conmutación online/offline}{red!70}
 \gbar{10}{10}{11}{T10 Pruebas y entrega}{red!70}
 \gbar{11.6}{1}{3}{Compra Lote 1 (núcleo)}{blue!45}
 \gbar{12.6}{2}{4}{Compra Lote 2 (AP + switch)}{blue!45}
 \gbar{13.6}{6}{8}{Compra Lote 3 (LTE + batería)}{blue!45}
 \node[anchor=west,font=\tiny] at (4,15) {\textcolor{red!70}{$\blacksquare$} actividad de la cadena crítica de recursos \quad \textcolor{blue!45}{$\blacksquare$} ventana de compra (segundo plano)};
\end{tikzpicture}
\caption{Cronograma de 11 semanas, ruta crítica y ventanas de compra escalonada.}
\label{fig:gantt}
\end{figure}

\textbf{Ruta crítica.} Por precedencias lógicas, sin restricción de recursos, la ruta crítica es T1 $\to$ T2 $\to$ T3 $\to$ T4 $\to$ T7 $\to$ T8 $\to$ T10, de 8.5 semanas; T5, T6 y T9 tienen 1 semana de holgura. Como un único integrante no puede ejecutar tareas en paralelo, las actividades se serializan (nivelación de recursos) y el plazo sube a 11 semanas; en esa cadena crítica de recursos (Fig.~\ref{fig:gantt}) todas las tareas tienen holgura nula, y T5, T6 y T9 son las primeras candidatas a reordenarse o recortarse. Los puntos de mayor riesgo son T2 y T4/T7: la configuración del almacenamiento y de la caché condiciona el contenido servido y la conmutación offline, y los retrasos de compra bloquean T2 y T3. Los riesgos de energía y conectividad (R1, R2) se asocian a las tareas que afectan en la Tabla~\ref{tab:risk}.

\subsection{Hitos por fase}
\begin{itemize}
  \item \textbf{Fase 1 (sem. 1--3):} informe de arquitectura, diagramas L2/L3, plan de compras, matriz de riesgos.
  \item \textbf{Fase 2 (sem. 4--7):} SO y almacenamiento (T2), red base (T3), Squid con contenido estático (T4), portal cautivo básico (T5).
  \item \textbf{Fase 3 (sem. 8--11):} plataforma en contenedores (T6), multimedia y contenidos (T7), dashboard y conmutación (T8--T9), pruebas y entrega (T10).
\end{itemize}

\subsection{Gestión del alcance}
Como la cadena crítica de recursos no tiene holgura, ante retrasos se recorta por prioridad:
\begin{itemize}
  \item \textbf{Esencial:} red L2/L3, DHCP/DNS, portal con perfiles, Squid y espejos, plataforma Spring Boot, modo offline.
  \item \textbf{Deseable:} Jellyfin con pre-transcodificación, dashboard con búsqueda.
  \item \textbf{Opcional:} SSL-Bump para equipos administrados, segundo AP, monitoreo avanzado.
\end{itemize}

\subsection{Recursos y plan de compras escalonado}
Los dispositivos se adquieren de forma progresiva, cuando cada fase los necesita, lo que distribuye la inversión (Tabla~\ref{tab:bom}). Antes de la llegada de cada lote se trabaja con alternativas: el desarrollo de la plataforma y las imágenes \texttt{arm64} se prueban en una laptop mediante Buildx/QEMU; hasta que llegue el router LTE (Lote~3) la WAN se toma de Ethernet del campus o del hogar.

\begin{table}[htbp]
\caption{Plan de compras escalonado y costo de referencia (verificar precios locales; la memoria DRAM/NAND tuvo alzas en 2026)}
\label{tab:bom}
\centering\footnotesize
\begin{tabularx}{\columnwidth}{l Y r l}
\toprule
\textbf{Lote} & \textbf{Ítems} & \textbf{USD} & \textbf{Pedido} \\
\midrule
1 & RPi5 8 GB, Active Cooler, carcasa, fuente 27 W; M.2 HAT+ con NVMe 256--512 GB; SSD USB 3 de 2 TB; adaptador USB 3 Gigabit & 380 & sem. 2 \\
2 & AP Wi-Fi 6 con VLAN y múltiples SSID; switch gestionable de 5 puertos & 110 & sem. 3 \\
3 & Router LTE con antenas externas y SIM de datos; batería LiFePO$_4$ 12 V/20 Ah con cargador DC-DC; conversor 12 V a 5 V & 200 & sem. 7 \\
Subtotal & Lotes 1--3 & 690 & \\
Contingencia & 15\,\% por variación de precios de memoria y almacenamiento, envío e impuestos & 104 & \\
\textbf{Total} & & \textbf{794} & \\
\bottomrule
\end{tabularx}
\end{table}

%=====================================================================
\section{Gestión de Riesgos}
La Tabla~\ref{tab:risk} valora probabilidad (P) e impacto (I) de 1 a 5; el nivel es $P\times I$.

\begin{table}[t]
\caption{Matriz de riesgos técnicos, ambientales y de gestión}
\label{tab:risk}
\centering\footnotesize\setlength{\tabcolsep}{3pt}
\begin{tabularx}{\textwidth}{l Y c c c l Y}
\toprule
\textbf{ID} & \textbf{Riesgo} & \textbf{P} & \textbf{I} & \textbf{Nivel} & \textbf{Afecta} & \textbf{Mitigación / plan de contingencia} \\
\midrule
R1 & Corte eléctrico en la unidad móvil & 4 & 5 & 20 & T2, T10 & Batería auxiliar LiFePO$_4$ + DC-DC, apagado ordenado por umbral de voltaje, \texttt{ext4} con \texttt{data=ordered}, \textit{fsck} automático. \\
R2 & Falta total de enlace satelital/móvil & 5 & 4 & 20 & T7, T9 & Modo offline (\texttt{offline\_mode}), DNS local, siembra inicial del contenido, sincronización diferida. \\
R3 & Sobrecarga de trabajo con un único integrante & 4 & 4 & 16 & Todas & Secuencia nivelada, entregas semanales incrementales, alcance recortable (Esencial/Deseable/Opcional). \\
R4 & Falla del SSD o corrupción del sistema de archivos & 3 & 5 & 15 & T2, T7 & SSD (resiste vibración), respaldos de BD y configuración, SMART, particiones separadas, imagen del sistema y respaldo en laptop (SSD de repuesto fuera del presupuesto). \\
R5 & Retraso o indisponibilidad de hardware por compra escalonada & 3 & 4 & 12 & T2, T3, T9 & Pedir con una semana de margen; desarrollo previo en laptop (Buildx arm64); proveedores locales alternos. \\
R6 & Sobrecalentamiento de la RPi5 en el vehículo & 4 & 3 & 12 & T10 & Active Cooler, carcasa ventilada, monitoreo de temperatura. \\
R7 & Saturación del SSD de 2 TB & 3 & 4 & 12 & T4, T7 & Particiones con cuotas, purga LFUDA, \textit{cache\_swap} 90/95\,\%, alertas al 85\,\%. \\
R8 & Exposición de datos de salud & 2 & 5 & 10 & T5, T6 & SSID WPA3 dedicado, VLAN aislada, HTTPS interno, roles y registros de acceso. \\
R9 & Portal cautivo incompatible con algún dispositivo & 3 & 3 & 9 & T5 & Pruebas con varios modelos; página de respaldo \texttt{portal.internal}; DHCP opción 114 si aplica. \\
R10 & Rendimiento insuficiente de la RPi5 (CPU/RAM/USB) & 3 & 3 & 9 & T7, T10 & Presupuesto de RAM (Tabla~\ref{tab:ram}), pre-transcodificar, mover caché pequeña a NVMe, pruebas de carga; criterio de migración a Mini PC. \\
R11 & Agotamiento del plan de datos LTE & 3 & 3 & 9 & T9 & Cuota de sincronización de 1~GB/día, límites \texttt{tc} por perfil, sincronización en horas de baja demanda, alerta de consumo. \\
R12 & Incumplimiento de protección de datos personales & 2 & 4 & 8 & T5, T6 & Registros mínimos con retención de 30 días, aviso de tratamiento en el portal, respaldos cifrados. \\
\bottomrule
\end{tabularx}
\end{table}

%=====================================================================
\section{Conclusiones y Trabajo Futuro}
El plan establece una arquitectura basada en una Raspberry Pi 5 que concentra enrutamiento, DNS/DHCP, portal cautivo, proxy/caché y servicios de contenido, con la plataforma Spring Boot en contenedores. Incluye segmentación L2/L3, un portal único con perfiles para los tres grupos, un SSID cifrado para equipos clínicos, y una estrategia de caché que combina siembra inicial, espejos locales y operación offline. La ruta crítica, el plan de compras escalonado y la gestión de alcance responden a la restricción de un único integrante y a la inversión progresiva. En la Fase~2 se construirá el prototipo (SO, red base, Squid y portal básico) y en la Fase~3 se integrarán la plataforma, el \textit{streaming}, el \textit{dashboard} y la conmutación automática.

\begin{thebibliography}{00}
\bibitem{rpi5} Raspberry Pi Ltd, ``Raspberry Pi 5 Product Brief,'' \url{https://datasheets.raspberrypi.com/rpi5/raspberry-pi-5-product-brief.pdf} (accedido el 16 oct. 2026).
\bibitem{rpiprice} OMG! Ubuntu, ``Raspberry Pi announces more price rises,'' feb. 2026. \url{https://www.omgubuntu.co.uk/2026/02/raspberry-pi-annouce-more-price-rises}
\bibitem{ieee8021q} IEEE, ``IEEE Standard for Local and Metropolitan Area Networks -- Bridges and Bridged Networks,'' IEEE Std 802.1Q-2018.
\bibitem{rfc1918} Y. Rekhter \textit{et al.}, ``Address Allocation for Private Internets,'' BCP 5, RFC 1918, Feb. 1996.
\bibitem{nft} Netfilter Project, ``nftables wiki,'' \url{https://wiki.nftables.org/} (accedido el 16 oct. 2026).
\bibitem{dnsmasq} S. Kelley, ``dnsmasq documentation,'' \url{https://thekelleys.org.uk/dnsmasq/doc.html} (accedido el 16 oct. 2026).
\bibitem{opennds} openNDS Documentation, ``Captive Portal Engine,'' \url{https://opennds.readthedocs.io/} (accedido el 16 oct. 2026).
\bibitem{squid} Squid Web Proxy Cache, ``Configuration Reference,'' \url{http://www.squid-cache.org/Doc/config/} (accedido el 16 oct. 2026).
\bibitem{nginx} Nginx, ``Documentation,'' \url{https://nginx.org/en/docs/} (accedido el 16 oct. 2026).
\bibitem{springboot} Spring, ``Spring Boot Reference Documentation,'' \url{https://docs.spring.io/spring-boot/} (accedido el 16 oct. 2026).
\bibitem{buildx} Docker, ``Multi-platform builds,'' \url{https://docs.docker.com/build/building/multi-platform/} (accedido el 16 oct. 2026).
\bibitem{postgres} PostgreSQL Global Development Group, ``PostgreSQL Documentation,'' \url{https://www.postgresql.org/docs/} (accedido el 16 oct. 2026).
\bibitem{compose} Docker, ``Docker Compose overview,'' \url{https://docs.docker.com/compose/} (accedido el 16 oct. 2026).
\bibitem{rfc8910} W. Kumari and E. Kline, ``Captive-Portal Identification in DHCP and Router Advertisements,'' RFC 8910, Sep. 2020.
\bibitem{internaltld} J. Davies, ``Private use top-level domain,'' IETF Internet-Draft draft-davies-internal-tld, ago. 2024.
\bibitem{rfc6762} S. Cheshire and M. Krochmal, ``Multicast DNS,'' RFC 6762, Feb. 2013.
\bibitem{kiwix} Kiwix, ``Offline Internet content,'' \url{https://kiwix.org/} (accedido el 16 oct. 2026).
\bibitem{kolibri} Learning Equality, ``Kolibri,'' \url{https://learningequality.org/kolibri/} (accedido el 16 oct. 2026).
\bibitem{jellyfin} Jellyfin, ``Documentation,'' \url{https://jellyfin.org/docs/} (accedido el 16 oct. 2026).
\bibitem{sslpeek} Squid Wiki, ``SslPeekAndSplice,'' \url{https://wiki.squid-cache.org/Features/SslPeekAndSplice} (accedido el 16 oct. 2026).
\bibitem{rpipower} Raspberry Pi Ltd, ``USB Power Delivery on Raspberry Pi 5,'' White Paper RP-009856.
\bibitem{ley1581} Congreso de la República de Colombia, ``Ley Estatutaria 1581 de 2012: disposiciones generales para la protección de datos personales,'' 2012.
\end{thebibliography}

\end{document}